%==============================================================================
% Combined Project 2 -- draft report
%
% DRAFT. Not submittable as-is. Every \NEEDSTEAM{...} below marks something only
% the team can supply (names, student IDs, individual contributions, the signed
% AI declaration). Search for NEEDSTEAM before submitting.
%
% Template: ACM Conference Proceedings Primary Article Template (acmart, sigconf).
% Drop this file and refs.bib into the Overleaf template and compile.
%
% PAGE BUDGET (6 pages excluding references, contribution text and appendices;
% 5 percentage points penalty per additional half-page):
%   Abstract 250-300 words + AI declaration   -- before Introduction, as required
%   Introduction + Literature Review          -- 1.5 pages
%   Data exploration and pre-processing       -- 0.5 page
%   Model development + evaluation            -- 1.0 page
%   XAI + discussion                          -- 1.5 pages
%   Emerging trends + ethics                  -- 1.0 page
%   Summary and conclusion                    -- 0.5 page
%
% This draft is written slightly under budget so the team's own interpretation
% can be added without overrunning. Check the page count after every edit.
%
% Every number in this document comes from RESULTS.md, which is generated from
% the JSON/CSV artefacts by make_report_tables.py. If a number here disagrees
% with RESULTS.md, RESULTS.md is right.
%==============================================================================
\documentclass[sigconf,nonacm]{acmart}

% Deliberately no \textcolor here: acmart loads xcolor itself and an explicit
% \usepackage can clash with its options. Bold is enough to find these.
\newcommand{\NEEDSTEAM}[1]{\textbf{[TEAM: #1]}}

\begin{document}

\title{Explainable, Edge-Efficient Olive Leaf Disease Classification}

%------------------------------------------------------------------------------
% NEEDSTEAM: one \author block per student, with student number in \affiliation.
%------------------------------------------------------------------------------
\author{\NEEDSTEAM{author names}}
\affiliation{%
  \institution{TU Dublin, Tallaght Campus}
  \city{Dublin}
  \country{Ireland}}

\begin{abstract}
Olive cultivation in the Mediterranean faces increasing pressure from foliar
disease, and automated leaf-level diagnosis is attractive precisely where expert
agronomists are scarce --- which is also where computing hardware is
constrained. This work asks two questions of a three-class olive leaf dataset
(healthy, olive peacock spot, \emph{Aculus olearius}): how ImageNet transfer
learning compares with training from scratch and with intrinsically interpretable
models, and how far the resulting model can be compressed before either its
accuracy or its explanations degrade.

Reproducing the published dataset split revealed train/test leakage that had to
be corrected first. Perceptual hashing found 119 duplicate pairs across the
boundary, of which byte-identical hashing found only 24; together with
burst-linked photographs, 359 training images were removed while the official
test split was left intact. All subsequent results use this corrected split.

On the corrected data, DenseNet121 and MobileNetV2 reach 0.939 and 0.930 test
macro-F1. A CNN trained from scratch reaches 0.871, and reproducing the
architecture originally chosen for this dataset reaches 0.810 --- against 0.401
reported before the correction, showing that figure was a product of leakage and
architecture jointly rather than evidence that scratch training fails here. The
best intrinsically interpretable model, logistic regression over 47 engineered
features, reaches 0.742 on the same test images, quantifying a 20-point cost of
interpretability.

Compression is effective but asymmetric. Full-integer quantisation makes
MobileNetV2 4.9$\times$ smaller and 5.6$\times$ faster for 0.005 macro-F1;
unstructured pruning reduces no latency at all and shows its saving only after
entropy coding. LIME explanations degrade in step with accuracy rather than
independently of it.
\end{abstract}

\maketitle

%------------------------------------------------------------------------------
% AI DECLARATION -- required position: after the Abstract, before the
% Introduction. Level 2 (AI-assisted idea generation and brainstorming); no AI
% content permitted in the final submission.
%
% NEEDSTEAM: this must be written and signed by the team. The paragraph below is
% a structural placeholder showing what has to be covered, NOT submittable text.
% draft-ethics.md section 8 records the team's decision and the verifiable facts
% it has to be consistent with (all commits on develop/model-compression carry
% Co-Authored-By trailers recording AI assistance).
%------------------------------------------------------------------------------
\section*{Declaration of AI Usage}
\NEEDSTEAM{Signed declaration, written by the team. It must state the scope of AI
assistance truthfully, confirm adherence to Level 2, and confirm that all
submitted text was written, verified and is owned by the authors. It must be
consistent with the repository's commit history, which records AI assistance
throughout. Do not submit this placeholder.}

%==============================================================================
\section{Introduction and Rationale}
%==============================================================================

Olive (\emph{Olea europaea} L.) is a defining crop of the Mediterranean basin.
Two foliar conditions dominate its disease burden: olive peacock spot
(\emph{Venturia oleaginea}), a fungal infection producing characteristic ringed
lesions, and infestation by the eriophyid mite \emph{Aculus olearius}, which
produces silvering and deformation rather than discrete lesions. Distinguishing
them matters because they call for different treatments, and misdiagnosis wastes
both chemical input and the grower's time.

\NEEDSTEAM{One to two sentences on the agronomic and climate motivation. The
proposal's figures --- Mediterranean warming 20\% faster than the global average,
yield losses of approximately 20\% --- are currently uncited. Either cite them or
state the motivation qualitatively. See draft-literature-review.md section 7.}

The dataset is the Olive Leaf Image Dataset, published on Kaggle under CC0, with
2{,}720 training and 680 test photographs across three balanced-ish classes. Two
properties made it a good choice for this project and one made it a demanding
one. It is single-leaf imagery against a plain background, so segmentation and
engineered features are both tractable; it is small enough to train repeatedly on
CPU; and, as Section~\ref{sec:data} shows, its published split contains leakage,
which turned data auditing from a formality into a substantive part of the work.

Two research questions follow, matching the two modules:

\begin{itemize}
\item \textbf{RQ~a.} How effectively do ImageNet-pretrained CNNs transfer to
  olive leaf disease classification compared with training from scratch and with
  an interpretable glass-box baseline, and what do XAI techniques reveal about
  them?
\item \textbf{RQ~b.} To what extent can post-training quantisation and pruning
  reduce model size and inference latency while preserving accuracy and
  explanation fidelity, enabling deployment on low-cost edge hardware?
\end{itemize}

%==============================================================================
\section{Literature Review}
%==============================================================================

\NEEDSTEAM{Condense draft-literature-review.md to roughly 900 words. It currently
runs much longer. The structure below is what the rubric rewards --- emerging
tools in the context of the application --- and the citation gaps flagged in its
Section 7 must be closed or the corresponding claims weakened before submission.}

\paragraph{Plant disease classification with CNNs.}
Mohanty et al.~\cite{mohanty2016} established the transfer-learning baseline for
leaf disease classification on PlantVillage, and the pattern --- an ImageNet
backbone with a replaced head --- remains the default. \NEEDSTEAM{Add the
olive-specific transfer-learning studies the proposal asserts exist, or restrict
this paragraph to the general plant-disease literature.}

\paragraph{Efficient architectures and model compression.}
MobileNetV2~\cite{sandler2018} makes depthwise-separable convolution the basis of
an architecture designed for constrained inference. Compression of already-trained
networks divides into pruning and quantisation. Han et
al.~\cite{han2016} combined pruning, trained quantisation and Huffman coding for
35--49$\times$ storage reduction, framing the problem as one of storage rather
than latency. Kuzmin et al.~\cite{kuzmin2023} compared the two families directly
across nine models and four tasks and found quantisation usually preferable,
with pruning competitive only at extreme ratios --- a finding
Section~\ref{sec:compression} reproduces on this dataset. Silva and
Almeida~\cite{silva2024} deploy compressed leaf-disease models to a Raspberry
Pi~4B and report 1.48--2.13$\times$ speed-ups, though on thermal rather than RGB
imagery.

\paragraph{Explainability.}
Grad-CAM~\cite{selvaraju2017} and its Grad-CAM++ refinement attribute a
convolutional prediction to spatial regions via gradients.
LIME~\cite{ribeiro2016} instead fits a sparse linear surrogate around a
perturbed neighbourhood of one input, requiring only a function from inputs to
class probabilities. That distinction is what makes LIME usable in
Section~\ref{sec:compression}: a quantised TFLite graph exposes no gradients.

%==============================================================================
\section{Data Exploration and Pre-processing}
\label{sec:data}
%==============================================================================

\paragraph{The published split leaks.}
Comparing 16$\times$16 greyscale thumbnails at a mean-absolute-error threshold
identified 119 duplicate photograph pairs spanning the train/test boundary. MD5
hashing finds only 24 of them: the remaining 95 are the same photograph
re-encoded, identical to the eye and to a downsampled comparison but not
byte-for-byte. A second class of leakage comes from burst photography, where
consecutive frames of one leaf are split across the boundary. Removing training
images that duplicate a test image (114) or share a burst with one (264) removes
359 images in total, leaving 2{,}361. The official test split was deliberately
left untouched so that results remain comparable with other work.

The effect is measurable. Before correction, the project's CNN reported
validation macro-F1 0.819 against test 0.946 --- a model scoring markedly higher
on data it was never validated against, which is the signature of test
contamination. After correction, every transfer-learned model shows the expected
ordering.

\paragraph{Capture source is confounded with label.}
Filenames encode the capture device. Cross-tabulating the prefix against the
class shows severe entanglement: the bare-numeric block is 100\%
\emph{Aculus olearius} and the \texttt{B*} block is 99.8\% healthy. A classifier
given \emph{nothing but the filename prefix} reaches 91.8\% accuracy on training
data against a 38.2\% majority baseline. It reaches only 49.9\% on the official
test split, which is itself evidence that the train and test photographs were
captured under different conditions. Any accuracy figure on this dataset must be
read with that confound in mind, and Section~\ref{sec:xai} tests whether the
model actually exploits it.

\paragraph{Segmentation.}
Eleven segmentation configurations across five model families were benchmarked on
the CVPPP A1 split; fine-tuned YOLO11-seg won on both quality (SBD 0.847, FBD
0.958) and speed (0.015\,s/image). Applied to the olive photographs it produced
3{,}003 files, which reconcile as 2{,}677 detected leaves plus 326 whole-image
fallbacks where no leaf was found. About 11\% of the crops feeding the glass-box
features are therefore unsegmented photographs with their background intact.

\paragraph{Pre-processing.}
All deep models share one pipeline: 224$\times$224 inputs, backbone-specific
ImageNet normalisation, training-only augmentation (rotation, zoom, shift,
flips), and balanced class weights. Holding this constant is what allows
architecture and compression effects to be separated from pre-processing effects.

%==============================================================================
\section{Model Development and Evaluation}
%==============================================================================

Four deep models were trained on the corrected split, differing only in
initialisation and architecture. Two are ImageNet-pretrained backbones fine-tuned
in two stages --- a head trained on a frozen backbone, then the final block
unfrozen at a lower learning rate with BatchNorm held in inference mode. Two are
randomly initialised: one reproducing layer-for-layer the architecture originally
selected for this dataset, and one redesigned to spend its parameters
differently. Model selection is by validation loss throughout, with the test
split evaluated once at the end.

Two interpretable models --- a depth-3 decision tree and logistic regression ---
were fitted to nine engineered features: five colour fractions measured inside a
leaf mask (green, yellow as a chlorosis proxy, brown, dark as necrosis, grey as
silvering) and four GLCM texture descriptors. Critically, these are evaluated
under the \emph{same} protocol as the CNNs, on the same official test images, so
the comparison in Table~\ref{tab:rqa} is like-for-like.

\begin{table}[t]
\caption{RQ~a: all models on the official 680-image test split.}
\label{tab:rqa}
\begin{tabular}{lrrl}
\toprule
Model & Params & macro-F1 & Explanation \\
\midrule
DenseNet121 (ImageNet)   & 7{,}040{,}579 & 0.939 & post-hoc \\
MobileNetV2 (ImageNet)   & 2{,}261{,}827 & 0.930 & post-hoc \\
Scratch CNN, redesigned  &   585{,}059 & 0.871 & post-hoc \\
Scratch CNN, original    & 5{,}767{,}139 & 0.810 & post-hoc \\
Logistic regression      &  47 features & 0.742 & intrinsic \\
Decision tree            &   9 features & 0.667 & intrinsic \\
\bottomrule
\end{tabular}
\end{table}

Three results follow, and only the first is the expected one.

Transfer learning wins, but by six points of macro-F1 over a competently designed
scratch network rather than by the margin the earlier work implied. What ImageNet
initialisation buys most clearly here is training time: MobileNetV2 reached 0.930
in 1{,}175\,s against 2{,}632\,s for the scratch network's 0.871.

The 0.401 macro-F1 previously reported for a scratch CNN on this dataset was not
evidence that scratch training fails. Reproducing that architecture unchanged on
the corrected split yields 0.810. The remaining gap to 0.871 is architectural:
the original design commits 5.75M of its 5.77M parameters to a single
\texttt{Flatten} into \texttt{Dense(64)}, so redirecting a tenth of that budget
into depth, batch normalisation and global average pooling beats it outright.
Both causes must be named; attributing the figure to either alone would be wrong.

One anomaly deserves reporting rather than smoothing. The redesigned scratch
network is the only model whose validation score (0.916) \emph{overstates} its
test score (0.871). Validation is carved from the training photographs while the
test split is the dataset's own, so a model whose features are learned entirely
from these photographs can lean on what they share --- consistent with the
capture-source confound above. Both transfer-learned models show the opposite,
smaller gap of $+0.015$.

%==============================================================================
\section{Explainability}
\label{sec:xai}
%==============================================================================

Explanations here serve four distinct purposes, and separating them is what keeps
the section from being decorative: debugging, bias detection, trust calibration,
and drift monitoring.

\paragraph{Debugging.} The first finding was negative. The saliency maps in the
project's earlier notebook explain the scratch CNN with 0.401 macro-F1 and 0.5\%
recall on \emph{Aculus olearius}, not the working model --- the wrong checkpoint
had been saved and reloaded. Every heatmap was faithfully explaining a model that
essentially never predicted one of the three classes.

\paragraph{Bias detection.} Grad-CAM++ attention was measured against a leaf
mask. The model places 10.9\% of its attention mass outside a leaf occupying
15.5\% of the frame --- a ratio of 0.70, meaning attention concentrates on leaf
tissue well beyond chance. The ratio is essentially unchanged between correct
(0.71) and incorrect (0.67) predictions. The capture-source confound is
demonstrably present in the data, and this model is demonstrably not exploiting
it. Establishing that required measurement; inspecting heatmaps by eye could not
have settled it.

\paragraph{Intrinsic against post-hoc.} The glass-box models offer something the
CNNs cannot: an explanation that \emph{is} the model. For logistic regression the
class score is a sum, so each feature's contribution is exactly
$\text{coefficient} \times \text{standardised value}$ with no surrogate involved.
For the depth-3 tree the explanation is the path taken, three tests with explicit
thresholds. Examining two correct and two incorrect test predictions per model
showed that both of the tree's errors turn on its very first split
(\texttt{frac\_yellow} $\leq 0.0134$), one of them missing the threshold by 0.03
standard deviations. The errors are not diffuse; they are a knife-edge on a
single feature. No post-hoc method could establish that.

The interpretability cost is 20 points of macro-F1 (0.939 against 0.742). Two
qualifications travel with that figure. It must be quoted from the common
protocol: the same models on segmented crops from training photographs reach
0.831, and pairing that with 0.939 would compare two different test sets. And it
depends on the feature set --- the nine base features reach only 0.667, while 47
features carry logistic regression to 0.742 and simultaneously cost the depth-3
tree, which can consult only three features however many it is offered. No claim
about which glass-box model is stronger survives without naming both the protocol
and the feature set.

%==============================================================================
\section{Emerging Techniques: Compression}
\label{sec:compression}
%==============================================================================

Two focus areas were implemented: transfer learning (Section~4) and model
compression, the latter in two variants --- post-training quantisation and
magnitude pruning.

\begin{table}[t]
\caption{MobileNetV2 compression. Latency is median single-image, single-thread
CPU.}
\label{tab:compression}
\begin{tabular}{lrrrr}
\toprule
Variant & MB & gzip MB & ms & macro-F1 \\
\midrule
Baseline (Keras)  & 13.21 & 8.74 & 53.55 & 0.930 \\
Float32 TFLite    &  8.88 & 8.23 & 11.56 & 0.930 \\
Float16           &  4.48 & 4.11 & 11.16 & 0.928 \\
Dynamic range     &  2.51 & 2.16 & 18.63 & 0.934 \\
Full integer      &  2.72 & 2.21 &  9.48 & 0.924 \\
Pruned 50\%       & 13.24 & 5.55 & 52.10 & 0.833 \\
\bottomrule
\end{tabular}
\end{table}

\paragraph{Quantisation is close to free.} Full-integer quantisation makes
MobileNetV2 4.9$\times$ smaller and 5.6$\times$ faster for 0.005 macro-F1. The
smallest model is not the fastest: dynamic-range quantisation produces the
smallest file but runs slower than float16, because its int8 weights are
dequantised at runtime.

\paragraph{Unstructured pruning does not accelerate anything.} Zeroing weights
without changing tensor shapes leaves the file size and the latency unchanged
--- 52.10\,ms against a 53.55\,ms baseline is measurement noise. The saving
appears only after entropy coding, where the gzipped size falls from 8.74\,MB to
5.55\,MB. Reporting both columns is what makes the result readable rather than
disappointing.

\paragraph{Architecture decides what compression costs.} MobileNetV2 loses
0.097 macro-F1 at 50\% sparsity; DenseNet121 at the same sparsity \emph{gains}
slightly (0.939 to 0.947), and tolerates 65\% before degrading. DenseNet121's
redundancy lies in the number of its weights, MobileNetV2's in their precision,
which is precisely why the architecture already designed for efficiency has less
to give up. This reproduces Kuzmin et al.~\cite{kuzmin2023} on our own data.

\paragraph{Explanation fidelity tracks accuracy.} LIME explanations were compared
before and after quantisation using a fixed segmentation and seed, so any
difference is attributable to the models. Rank correlation of the feature weights
runs 0.599 (float16), 0.556 (dynamic range) and 0.407 (full integer) --- the
same ordering as the accuracy cost. Top-5 region overlap is higher than the rank
correlation, locating the instability in low-weight regions rather than in the
evidence a prediction rests on. Label agreement was 6/6 images for float16 and
full integer, and 5/6 for dynamic range; six images cannot support a rate, which
is why the raw counts are given. Fidelity does not degrade independently of
accuracy.

%==============================================================================
\section{Ethical Considerations}
%==============================================================================

Four positions were decided by the team, each grounded in a measurement rather
than in generic principle.

\paragraph{Triage, not diagnosis.} The system is framed as a triage and
decision-support tool that prioritises which trees an expert inspects. Two
measurements rule out a diagnostic claim: the label set is closed at three
classes, so any condition outside it is silently mapped onto one of them, and the
errors that occur are confident ones between diseases rather than between
diseased and healthy.

\paragraph{No raw confidence.} Confidence is not calibrated and the failures show
it: one misclassification carries 0.979 confidence and another 0.991. No raw
softmax output is presented to a grower as a probability; the system abstains
below a threshold. Displaying a numeric confidence requires calibration first,
which this project has not performed and which is stated as future work rather
than implied to be done.

\paragraph{Metadata.} EXIF is stripped by default from any image shared outward,
with provenance retained separately where the bias analysis needs it. The
pipeline already drops EXIF as a side effect of re-encoding; the decision makes
an incidental behaviour an intended one.

\paragraph{Provenance and licensing.} The dataset is CC0, so attribution is an
academic norm rather than a legal obligation, and it is given. Provenance is
incomplete --- the dataset documents neither who photographed the leaves nor how
labels were assigned --- and the leakage found during preparation is a reason for
caution about its curation generally.

\paragraph{Model drift.} Treating drift as a post-hoc XAI concern, the bias
measurements above form a monitoring baseline. The realistic drift axis is
capture conditions, not the disease: a deployment that changed camera, season or
lighting would move the input distribution along exactly the axis this dataset
confounds with the label. Three signals are monitorable in increasing cost ---
input colour and device statistics, prediction distribution and confidence, and
the attention-outside-leaf ratio, currently 0.70 and trending toward 1.0 if the
model started reading background. One measurement per signal sets a baseline but
cannot set a threshold; a sustained move across batches, sized against the
variation between existing capture sources, is the defensible trigger. Response
escalates with the signal: re-measure for an input-side change alone, retrain on
data including the new conditions if the attention ratio rises, and suspend
deployment if confidence collapses without an input-side explanation.

%==============================================================================
\section{Summary and Conclusion}
%==============================================================================

Auditing the published dataset split was not preliminary work but a precondition
for every number reported here: 119 leaked duplicate pairs, of which exact
hashing found 24, plus burst-linked photographs, required removing 359 training
images. The corrected split reversed the previously reported
validation-above-test inversion and cut a previously reported scratch-CNN result
approximately in half in the opposite direction --- 0.401 becoming 0.810 once
leakage and architecture were separated.

On RQ~a, ImageNet transfer reaches 0.939 macro-F1 against 0.871 from scratch and
0.742 from the best intrinsically interpretable model, all on the same test
split. The transfer advantage is real but smaller than expected, and expresses
itself as much in training time as in final accuracy. On RQ~b, quantisation
delivers a 4.9$\times$ size and 5.6$\times$ latency improvement for 0.005
macro-F1, while unstructured pruning improves neither size nor latency and is
architecture-dependent in what it costs. Explanation fidelity degrades in step
with accuracy, so compression does not silently trade away interpretability as a
separate axis.

\NEEDSTEAM{Add two to three sentences of the team's own conclusion: what you
would do next, and what you would tell a grower considering this system.}

Limitations are stated where they arise: LIME fidelity rests on six images,
Grad-CAM++ output at a 7$\times$7 grid cannot localise individual lesions, three
pruning sweep points did not converge, and the capture-source confound bounds
what any accuracy figure on this dataset means.

%==============================================================================
% Excluded from the 6-page limit: references, contributions, appendices.
%==============================================================================
\bibliographystyle{ACM-Reference-Format}
\bibliography{refs}

\appendix
\section{Individual Contributions}
\NEEDSTEAM{One subsection per student, unique text, maximum 500 words each. This
is excluded from the page limit and is graded individually. Each student uploads
the project and their own contribution.}

\end{document}
